% Propietario conservado: /Users/ruben/Documents/New project/output/EXTREMA_MEDIA_RAZON_RECIPROCIDAD_ES_EN_20260918/ARTICULO/ES/source/libro/rectas27.tex
\section{Holonomie finie, configuration de Schläfli et symétrie \texorpdfstring{\(E_6\)}{E6}}
\label{x_reciprocidad:p12-83-c20-s6:chap:holonomia-schlafli-e6-nuevo}
\index{holonomie finie}
\index{Schläfli}
\index{E6@\(E_6\)}

La branche de Witt n'épuise pas les structures exceptionnelles accessibles à partir d'APP. La branche parallèle de ce chapitre part des deux polarisations mixtes somme--produit et produit--somme, avec une orientation de plaquette et une section centrale déclarées. La courbure calculée détermine la forme alternée ; celle-ci définit l'extension centrale. Une extension n'est pas inférée de la seule cardinalité 27.

\subsection{Courbures mixtes d’APP}

Sur \(T_9=(\mathbb Z/9\mathbb Z)^2\), avec
\(S(x,y)=x+y\) et \(P(x,y)=xy\), soient
\[
 \Delta_xf=f(x+1,y)-f(x,y),\qquad
 \Delta_yf=f(x,y+1)-f(x,y).
\]
Les connexions mixtes sont
\[
 A^{SP}=(\Delta_xS,\Delta_yP)=(1,x),\qquad
 A^{PS}=(\Delta_xP,\Delta_yS)=(y,1).
\]
Leur courbure discrète est
\[
 F_A=\Delta_xA_y-\Delta_yA_x.
\]

\begin{proposition}[Courbures orientées opposées]
\label{x_reciprocidad:p12-83-c20-s6:prop:curvaturas-app-e6-nuevo}
\[
 F_{A^{SP}}=1,\qquad F_{A^{PS}}=-1\pmod9.
\]
La circulation sur un rectangle orienté de côtés \(r,s\) est,
respectivement, \(rs\) et \(-rs\).
\end{proposition}

\begin{proof}
\(\Delta_xx-\Delta_y1=1\) et
\(\Delta_x1-\Delta_yy=-1\). Un rectangle contient \(rs\) plaquettes ; lors de la
sommation, les arêtes intérieures s’annulent et il reste la circulation du bord.
\end{proof}

Pour des déplacements \(u=(a,b)\), \(v=(c,d)\), la classe alternée est
\begin{equation}
 \sigma(u,v)=ad-bc\pmod9.
 \label{x_reciprocidad:p12-83-c20-s6:eq:forma-alternante-e6-nueva}
\end{equation}
L’orientation fixe le signe de \(\sigma\). La section centrale employée
ci-dessous appartient à la jauge de cette réalisation ; la modifier par un cobord
produit une extension isomorphe, non un nouvel invariant.

\subsection{Extension centrale et réduction ternaire}

Nous définissons le groupe de Heisenberg nonadique
\begin{equation}
 \mathcal H_9=\{(a,b,z):a,b,z\in\mathbb Z/9\mathbb Z\},
 \quad
 (a,b,z)(c,d,w)=(a+c,b+d,z+w+bc).
 \label{x_reciprocidad:p12-83-c20-s6:eq:heisenberg-9-nuevo}
\end{equation}
Le commutateur de deux déplacements horizontaux est
\[
 [(a,b,0),(c,d,0)]=(0,0,bc-ad),
\]
qui enregistre \(-\sigma(u,v)\). En réduisant chaque coordonnée modulo trois, on
obtient
\begin{equation}
 H_3=\{(a,b,z):a,b,z\in\mathbb F_3\},\qquad |H_3|=27,
 \label{x_reciprocidad:p12-83-c20-s6:eq:heisenberg-3-nuevo}
\end{equation}
avec la même loi. Le noyau de \(\mathcal H_9\to H_3\) est \((3\mathbb Z/9\mathbb Z)^3\), de cardinal 27.

\subsection{Représentation de Weyl--Heisenberg et phase centrale}

L’extension nonadique admet une réalisation unitaire explicite. Soit
\[
 \zeta_9:=\exp(2\pi i/9),
 \qquad
 \mathscr H_9^{\rm Sch}:=\mathbb C^9,
\]
de base \(\{|n\rangle:n\in\mathbb Z/9\mathbb Z\}\). Nous définissons
\begin{equation}
 X|n\rangle:=|n+1\rangle,
 \qquad
 Z|n\rangle:=\zeta_9^n|n\rangle.
 \label{x_reciprocidad:p12-83-c20-s6:eq:weyl-desplazamiento-fase-u029}
\end{equation}
On a
\[
 ZX=\zeta_9XZ
\]
et, par multiplication,
\begin{equation}
 (X^aZ^b)(X^cZ^d)
 =\zeta_9^{bc}X^{a+c}Z^{b+d}.
 \label{x_reciprocidad:p12-83-c20-s6:eq:cociclo-weyl-nonadico-u029}
\end{equation}

\begin{theorem}[Représentation fidèle de Weyl--Heisenberg]
\label{x_reciprocidad:p12-83-c20-s6:thm:representacion-weyl-heisenberg-u029}
L'application
\begin{equation}
 \rho_{\rm WH}:\mathcal H_9\longrightarrow
 U(\mathscr H_9^{\rm Sch}),
 \qquad
 \rho_{\rm WH}(a,b,z):=\zeta_9^zX^aZ^b,
 \label{x_reciprocidad:p12-83-c20-s6:eq:representacion-weyl-heisenberg-u029}
\end{equation}
est une représentation unitaire fidèle. Pour
\(u=(a,b)\) et \(v=(c,d)\), son commutateur satisfait
\begin{equation}
 [\rho_{\rm WH}(u,0),\rho_{\rm WH}(v,0)]
 =\zeta_9^{\,bc-ad}I
 =\zeta_9^{-\sigma(u,v)}I.
 \label{x_reciprocidad:p12-83-c20-s6:eq:conmutador-weyl-holonomia-u029}
\end{equation}
\end{theorem}

\begin{proof}
L’identité \eqref{x_reciprocidad:p12-83-c20-s6:eq:cociclo-weyl-nonadico-u029} reproduit le terme
\(bc\) de \eqref{x_reciprocidad:p12-83-c20-s6:eq:heisenberg-9-nuevo} ; d’où la propriété
d’homomorphisme. Si \(\zeta_9^zX^aZ^b=I\), l’action sur la base impose
successivement \(a=0\), \(b=0\) et \(z=0\). Cela prouve la fidélité. Comparer
les produits dans les deux ordres donne
\(\zeta_9^{bc-ad}\), qui est la phase alternée de
\eqref{x_reciprocidad:p12-83-c20-s6:eq:forma-alternante-e6-nueva}.
\end{proof}

La représentation distingue deux observables conjugués. Soient
\[
 Q_{\rm H}:=\operatorname{diag}(0,1,\ldots,8),
 \qquad
 F_{jk}:=\frac13\zeta_9^{jk},
 \qquad
 P_{\rm H}:=F^\dagger Q_{\rm H}F.
\]
Les bases propres de \(Q_{\rm H}\) et \(P_{\rm H}\) sont mutuellement non biaisées, puisque \(|F_{jk}|^2=1/9\). Pour tout vecteur unitaire \(\psi\in\mathscr H_9^{\rm Sch}\),
\begin{align}
 \operatorname{Var}_\psi(Q_{\rm H})
 \operatorname{Var}_\psi(P_{\rm H})
 &\ge
 \frac14\left|
 \langle\psi,[Q_{\rm H},P_{\rm H}]\psi\rangle
 \right|^2,
 \label{x_reciprocidad:p12-83-c20-s6:eq:incertidumbre-weyl-u029}\\
 H_{Q_{\rm H}}(\psi)+H_{P_{\rm H}}(\psi)
 &\ge\log9.
 \label{x_reciprocidad:p12-83-c20-s6:eq:incertidumbre-entropica-weyl-u029}
\end{align}
Ici, \(H_{Q_{\rm H}}\) et \(H_{P_{\rm H}}\) sont les entropies de Shannon
des distributions de Born dans les bases propres respectives,
avec logarithme naturel et convention \(0\log0=0\).
Ces opérateurs sont adimensionnels. Une carte physique d’unités est une
réalisation ultérieure et n’intervient pas dans l’extension centrale.


\subsection{Action de 10 déplacements sur le graphe de Cayley}

Pour une forme linéaire \(\ell(a,b)=pa+qb\), nous posons
\[
 D_\ell(a,b)=\left(a,b,\frac{ab}{2}+\ell(a,b)\right),
 \qquad Z=(0,0,1),
\]
où \(2^{-1}=2\) dans \(\mathbb F_3\). L'ensemble stable par inversion
\begin{equation}
 S_\ell=
 \{D_\ell(v):v\in\mathbb F_3^2\setminus\{0\}\}
 \cup\{Z,Z^{-1}\}
 \label{x_reciprocidad:p12-83-c20-s6:eq:conjunto-conexion-e6}
\end{equation}
possède dix éléments et ne contient pas l'identité.

\begin{theorem}[Indépendance de la jauge linéaire]
\label{x_reciprocidad:p12-83-c20-s6:thm:independencia-calibre-e6}
Les neuf fonctionnelles \(\ell\) produisent des graphes de Cayley isomorphes. Les
neuf ensembles \(S_\ell\) forment une seule orbite sous
\(\operatorname{Aut}(H_3)\).
\end{theorem}

\begin{proof}
L'application
\[
 \phi_\ell(a,b,z)=(a,b,z+\ell(a,b))
\]
est un automorphisme central et satisfait
\(\phi_\ell(S_0)=S_\ell\). Pour écrire la famille complète dans les
coordonnées polarisées de \cref{x_reciprocidad:p12-83-c20-s6:eq:heisenberg-9-nuevo}, soit
\(c(v,w)=b c\) pour \(v=(a,b)\), \(w=(c,d)\). Étant donné
\(M\in\operatorname{GL}(2,3)\), le défaut

\[
 B_M(v,w)=c(Mv,Mw)-\det(M)c(v,w)
\]
est symétrique, car les parties alternées des deux termes coïncident. Nous définissons la correction quadratique
\begin{equation}
 q_M(v)=2B_M(v,v),
 \qquad
 q_M(v+w)-q_M(v)-q_M(w)=B_M(v,w),
 \label{x_reciprocidad:p12-83-c20-s6:eq:correccion-cuadratica-automorfismos-h3}
\end{equation}
où \(2=2^{-1}\) dans \(\mathbb F_3\). La famille complète d'automorphismes s'écrit alors
\begin{equation}
 (v,t)\longmapsto
 \bigl(Mv,\det(M)t+q_M(v)+\ell(v)\bigr),
 \qquad M\in\operatorname{GL}(2,3),\
 \ell\in(\mathbb F_3^2)^*.
 \label{x_reciprocidad:p12-83-c20-s6:eq:automorfismos-h3-completos}
\end{equation}
L’identité polaire de \cref{x_reciprocidad:p12-83-c20-s6:eq:correccion-cuadratica-automorfismos-h3}
prouve directement que le produit central est conservé. La famille possède
\(48\cdot9=432\) éléments. L’énumération exhaustive des
\(\binom{13}{5}=1287\) sous-ensembles stables par inversion de cardinal dix
trouve exactement ces neuf sous-ensembles avec les paramètres dérivés
ci-après.
\end{proof}

Écrivons \(\underline S=\sum_{s\in S}s\) dans
\(\mathbb Z[H_3]\). Le dénombrement des produits ordonnés donne
\begin{equation}
 \underline S^{\,2}
 =10e+\underline S+5(\underline{H_3}-e-\underline S)
 =5\underline{H_3}+5e-4\underline S.
 \label{x_reciprocidad:p12-83-c20-s6:eq:identidad-grupo-schlafli-nueva}
\end{equation}
En effet, l'identité reçoit dix factorisations ; un élément de \(S\) en reçoit une ; chacun des seize autres en reçoit cinq.

\begin{theorem}[Graphe des vingt-sept droites]
\label{x_reciprocidad:p12-83-c20-s6:thm:grafo-27-rectas-nuevo}
La matrice d’adjacence \(A_\ell\) de
\(\Gamma_\ell=\operatorname{Cay}(H_3,S_\ell)\) satisfait
\begin{equation}
 A_\ell^2=5I-4A_\ell+5J.
 \label{x_reciprocidad:p12-83-c20-s6:eq:identidad-schlafli-nueva}
\end{equation}
Par conséquent,
\[
 \Gamma_\ell=\operatorname{srg}(27,10,1,5),
 \qquad
 \operatorname{Spec}(A_\ell)=\{10^1,1^{20},(-5)^6\}.
\]
Il est isomorphe au graphe d'intersection des vingt-sept droites d'une surface cubique lisse.
\end{theorem}

\begin{proof}
La représentation régulière de
\cref{x_reciprocidad:p12-83-c20-s6:eq:identidad-grupo-schlafli-nueva} produit
\cref{x_reciprocidad:p12-83-c20-s6:eq:identidad-schlafli-nueva}. Ses entrées indiquent que tout sommet
possède dix voisins, qu’une paire adjacente possède un voisin commun et qu’une paire non
adjacente en possède cinq. Sur la droite constante, \(A_\ell\) agit par 10 ;
dans son complément, \(J=0\) et
\((A_\ell-I)(A_\ell+5I)=0\). La trace nulle fixe les multiplicités 20 et 6.

Dans le modèle d’une cubique obtenue en éclatant six points de
\(\mathbb P^2\), les 27 classes sont
\[
 E_i,\qquad L_{ij}=H-E_i-E_j,\qquad
 Q_i=2H-\sum_{j\ne i}E_j.
\]
Avec \(H^2=1\), \(E_i^2=-1\), \(H\cdot E_i=0\) et \(E_i\cdot E_j=0\) pour \(i\ne j\), la matrice d'adjacence définie par l'intersection entre classes distinctes possède les mêmes paramètres. La matrice complète d'intersection a pour diagonale \(-1\). L'identification à la configuration classique utilise le théorème d'unicité de cette réalisation des 27 droites \cite{x_reciprocidad:Manivel2005}, avec le domaine et l'action qui y sont spécifiés.
\end{proof}

L’assertion d’isomorphie peut être rendue entièrement effective par
une bijection calculée. L’étiquetage n’est pas absolu : une autre section linéaire
ou un automorphisme de la surface produit un autre tableau de la même classe.
\begin{longtable}{@{}c@{\,}c@{\,}c@{\qquad}l@{}}
\caption{Bijection explicite entre \(H_3\) et les vingt-sept classes de
droites.}
\label{x_reciprocidad:p12-83-c20-s6:tab:h3-27-rectas-u029}\\
\toprule
\(a\)&\(b\)&\(z\)&classe de droite\\
\midrule
\endfirsthead
\multicolumn{4}{@{}l}{\small\itshape Tableau \thetable\ (suite)}\\
\toprule
\(a\)&\(b\)&\(z\)&classe de droite\\
\midrule
\endhead
0&0&0&\(E_1\)\\
0&0&1&\(L_{12}\)\\
0&0&2&\(Q_2\)\\
0&1&0&\(L_{13}\)\\
0&1&1&\(Q_1\)\\
0&1&2&\(E_3\)\\
0&2&0&\(Q_3\)\\
0&2&1&\(E_2\)\\
0&2&2&\(L_{23}\)\\
1&0&0&\(L_{14}\)\\
1&0&1&\(L_{35}\)\\
1&0&2&\(L_{26}\)\\
1&1&0&\(L_{36}\)\\
1&1&1&\(L_{24}\)\\
1&1&2&\(L_{15}\)\\
1&2&0&\(L_{25}\)\\
1&2&1&\(L_{16}\)\\
1&2&2&\(L_{34}\)\\
2&0&0&\(Q_4\)\\
2&0&1&\(L_{46}\)\\
2&0&2&\(E_6\)\\
2&1&0&\(E_5\)\\
2&1&1&\(Q_6\)\\
2&1&2&\(L_{56}\)\\
2&2&0&\(L_{45}\)\\
2&2&1&\(E_4\)\\
2&2&2&\(Q_5\)\\
\bottomrule
\end{longtable}

Comme contrôle local, les dix voisins de \(e=(0,0,0)\) sont les éléments de
\(S_0\). Le tableau les transforme en les dix classes qui rencontrent \(E_1\) :
les cinq \(L_{1j}\) et les cinq \(Q_j\), avec \(2\le j\le6\). La
vérification complète compare les 729 entrées des deux matrices
d’adjacence.

Le complément \(B_\ell=J-I-A_\ell\) vérifie
\begin{equation}
 B_\ell^2=8I+2B_\ell+8J,
 \qquad
 \operatorname{srg}(27,16,10,8),
 \qquad
 \operatorname{Spec}(B_\ell)=\{16^1,4^6,(-2)^{20}\}.
 \label{x_reciprocidad:p12-83-c20-s6:eq:complemento-schlafli}
\end{equation}
C'est le graphe de Schläfli dans la convention d'adjacence complémentaire ; \(A_\ell\) est le graphe d'intersection de paramètres \((27,10,1,5)\).

Chaque arête appartient à un unique triangle, puisque \(\lambda=1\). Le nombre
de triangles est
\[
 \frac{27\cdot10}{6}=45.
\]
Le groupe complet des automorphismes est d’ordre
\[
 27\cdot1920=51840
\]
et, par son action sur la configuration de droites, il est identifié à \(W(E_6)\), et pas seulement par égalité des ordres \cite{x_reciprocidad:Manivel2005}.

\subsection{Le réseau de racines qui réalise \texorpdfstring{\(E_6\)}{E6}}

Pour expliciter l’objet sur lequel agit ce groupe, soit
\begin{equation}
 \operatorname{Pic}(X)
 =\mathbb ZH\oplus\bigoplus_{i=1}^6\mathbb ZE_i,
 \qquad
 H^2=1,\qquad E_i^2=-1,\qquad H\cdot E_i=E_i\cdot E_j=0\ (i\ne j),
 \label{x_reciprocidad:p12-83-c20-s6:eq:picard-cubica-e6}
\end{equation}
le réseau de Picard de l’éclatement de six points généraux de
\(\mathbb P^2\). Sa classe canonique est
\[
 K_X=-3H+E_1+\cdots+E_6.
\]
Le complément orthogonal
\begin{equation}
 L_{E_6}=K_X^\perp\subset\operatorname{Pic}(X),
 \qquad (x,y)_{E_6}:=-x\cdot y,
 \label{x_reciprocidad:p12-83-c20-s6:eq:reticulo-raices-e6}
\end{equation}
est défini positif de rang six. Une base de racines simples est
\begin{equation}
 \begin{aligned}
 \alpha_1&=E_1-E_2,& \alpha_2&=E_2-E_3,&
 \alpha_3&=E_3-E_4,\\
 \alpha_4&=E_4-E_5,& \alpha_5&=E_5-E_6,&
 \alpha_6&=H-E_1-E_2-E_3.
 \end{aligned}
 \label{x_reciprocidad:p12-83-c20-s6:eq:raices-simples-e6}
\end{equation}
Chaque \(\alpha_i\) est orthogonal à \(K_X\), est de norme deux et leurs produits
sont \(-1\) exactement pour les paires adjacentes du diagramme
\(E_6\) : la chaîne \(1-2-3-4-5\), avec le nœud \(6\) relié au nœud \(3\).
Par conséquent, leur matrice de Gram est la matrice de Cartan de \(E_6\).

\begin{proposition}[Les soixante-douze racines et leurs réflexions]
\label{x_reciprocidad:p12-83-c20-s6:prop:setenta-dos-raices-e6}
Les racines de \(L_{E_6}\) sont
\begin{equation}
 \begin{aligned}
 &E_i-E_j &&(i\ne j),\\
 &\pm(H-E_i-E_j-E_k) &&(1\le i<j<k\le6),\\
 &\pm(2H-E_1-\cdots-E_6),
 \end{aligned}
 \label{x_reciprocidad:p12-83-c20-s6:eq:enumeracion-raices-e6}
\end{equation}
au total \(30+40+2=72\). Pour chaque racine \(\alpha\),
\begin{equation}
 s_\alpha(x)=x+(x\cdot\alpha)\alpha
 \label{x_reciprocidad:p12-83-c20-s6:eq:reflexion-raiz-e6}
\end{equation}
préserve l’intersection et \(K_X\). Les six réflexions simples engendrent
\(W(E_6)\) et permutent fidèlement les 27 classes
\(E_i,L_{ij},Q_i\) de \cref{x_reciprocidad:p12-83-c20-s6:thm:grafo-27-rectas-nuevo}.
\end{proposition}

\begin{proof}
Substituer chaque classe de \cref{x_reciprocidad:p12-83-c20-s6:eq:enumeracion-raices-e6} dans \cref{x_reciprocidad:p12-83-c20-s6:eq:picard-cubica-e6} donne \(K_X\cdot\alpha=0\) et \(\alpha^2=-2\). Le décompte des classes est celui indiqué. La formule de réflexion est la réflexion orthogonale pour la forme \(-(\cdot)\) ; elle préserve \(K_X\) car \(K_X\cdot\alpha=0\). Le système simple de \cref{x_reciprocidad:p12-83-c20-s6:eq:raices-simples-e6} engendre toutes les racines et son groupe de Coxeter est \(W(E_6)\). Son action sur les 27 classes de courbes exceptionnelles conserve l'intersection, de sorte qu'elle coïncide avec l'action sur le graphe déjà construit \cite{x_reciprocidad:Manivel2005}.
\end{proof}

\subsection{Stratification des 729 mots}

Nous écrivons un mot de six trits sous la forme
\[
 w=(\ell_1,h_1,\ell_2,h_2,\ell_3,h_3),
\]
et nous définissons deux éléments de \(H_3\),
\[
 L(w)=(\ell_1,\ell_2,\ell_3),\qquad
 H(w)=(h_1,h_2,h_3),
\]
en utilisant la loi centrale dans ses produits. Selon
\(L(w)^{-1}H(w)\), les 729 mots se divisent en
\begin{align*}
 \mathcal D_\ell&=\{w:L(w)^{-1}H(w)=e\},\\
 \mathcal I_\ell&=\{w:L(w)^{-1}H(w)\in S_\ell\},\\
 \mathcal S_\ell&=\{w:L(w)^{-1}H(w)\notin\{e\}\cup S_\ell\}.
\end{align*}

\begin{theorem}[Partition Hensel--Schläfli]
\label{x_reciprocidad:p12-83-c20-s6:thm:particion-hensel-schlafli-nueva}
\begin{equation}
 729=27+270+432,
 \label{x_reciprocidad:p12-83-c20-s6:eq:particion-729-e6-nueva}
\end{equation}
correspondant à la diagonale, à l’incidence et au gauchissement.
\end{theorem}

\begin{proof}
Pour chacune des 27 valeurs de \(L(w)\), il y a une unique valeur de
\(H(w)\) dont la différence est l’identité, dix dont la différence est dans \(S_\ell\) et
seize autres. On multiplie \(27\) par \(1,10,16\).
\end{proof}

Le cardinal \(432\) de la strate gauchie coïncide numériquement avec la
multiplicité de la région transversale de \(\pi\) dans
\cref{x_reciprocidad:pi-estrella:figura}. Cela n’identifie pas les deux objets : ici,
on compte des mots classés par \(L(w)^{-1}H(w)\), tandis que là,
on compte des préimages d’un registre \(U_6\). Sans opérateur d’entrelacement explicite et
équivariant entre ces domaines, \(432=432\) n’est qu’un contrôle cardinal.

\subsection{Nécessité de la phase centrale}

L’énumération de tous les sous-ensembles stables par inversion de cardinal dix
dans les cinq groupes d’ordre 27 donne
\[
\begin{array}{c|c|c}
\text{groupe}&\text{abélien}&
\#\operatorname{srg}(27,10,1,5)\\ \hline
C_{27}&\text{oui}&0\\
C_9\times C_3&\text{oui}&0\\
C_3^3&\text{oui}&0\\
C_9\rtimes C_3&\text{non}&9\\
H_3&\text{non}&9.
\end{array}
\]
Ainsi, conserver 27 étiquettes mais abélianiser la composition détruit la configuration. Ni la cardinalité ni le degré ne suffisent.

Si \(\mathcal T\) est l’ensemble des 45 triangles, le cubique
\[
 \mathcal C_\Gamma(x_1,\ldots,x_{27})
 =\sum_{\{i,j,k\}\in\mathcal T}x_ix_jx_k
\]
récupère l'adjacence, car chaque arête appartient à un triangle unique. Son stabilisateur par permutations est exactement
\[
 \operatorname{Aut}(\Gamma_\ell)\cong W(E_6).
\]

Pour le diagramme suivant, on écrit
\(\Gamma_{27}=\Gamma_\ell\) et \(\mathcal T_{45}=\mathcal T\).
\begin{figure}[htbp]
\centering
\[
 \begin{gathered}
 \text{courbures APP }\pm1
 \longrightarrow\sigma
 \longrightarrow\mathcal H_9
 \longrightarrow H_3,\\[3pt]
 H_3
 \longrightarrow\Gamma_{27}
 \longrightarrow\mathcal T_{45}
 \longrightarrow W(E_6).
 \end{gathered}
\]
\caption{Chaîne causale de la branche d’holonomie mixte. Elle est parallèle à
\(C_W\to W_{12}\to M_{12}\) et n’en est pas un maillon.}
\label{x_reciprocidad:p12-83-c20-s6:fig:cadena-holonomia-e6-u029}
\end{figure}

\begin{criterion}[Critères de réfutation de la branche \(E_6\)]
\label{x_reciprocidad:p12-83-c20-s6:crit:refutacion-rama-e6-u029}
La branche est réfutée si les courbures ne sont pas opposées, si le cocycle
central est effacé, si la correction quadratique
\eqref{x_reciprocidad:p12-83-c20-s6:eq:correccion-cuadratica-automorfismos-h3} ne conserve pas le produit,
si \eqref{x_reciprocidad:p12-83-c20-s6:eq:identidad-grupo-schlafli-nueva} ne produit pas les coefficients
\((10,1,5)\), si un groupe abélien produit la même configuration, si la
partition n’a pas pour somme 729, si le tableau
\ref{x_reciprocidad:p12-83-c20-s6:tab:h3-27-rectas-u029} n’entrelace pas les matrices d’adjacence ou si
\(W(E_6)\) est identifié uniquement par son ordre et non par son action sur les
vingt-sept classes et les 72 racines.
\end{criterion}
